(1) 总势能为 \(U=\frac12k_0(x_1^2+x_2^2)+\frac12k(x_1-x_2)^2\)，动能为 \(T=m(\dot x_1^2+\dot x_2^2)/2\)。Euler–Lagrange 方程为 \(m\ddot x_1+(k_0+k)x_1-kx_2=0\)、\(m\ddot x_2+(k_0+k)x_2-kx_1=0\)，刚度矩阵对称。
(2) 用 \(q_+=(x_1+x_2)/\sqrt2\)、\(q_-=(x_1-x_2)/\sqrt2\) 正交变换，运动方程解耦为 \(\ddot q_++(k_0/m)q_+=0\)、\(\ddot q_-+[(k_0+2k)/m]q_-=0\)。对称模式形状 \((1,1)/\sqrt2\) 对应 \(\omega_+=\sqrt{k_0/m}\)；反对称模式 \((1,-1)/\sqrt2\) 对应 \(\omega_-=\sqrt{(k_0+2k)/m}\)，两形状内积为零。
(3) 初值使 \(q_+(0)=q_-(0)=A/\sqrt2\)，因此 \(x_1(t)=\frac A2[\cos(\omega_+t)+\cos(\omega_-t)]\)、\(x_2(t)=\frac A2[\cos(\omega_+t)-\cos(\omega_-t)]\)。利用和差化积，\(x_1=A\cos[(\omega_-+\omega_+)t/2]\cos[(\omega_--\omega_+)t/2]\)；弱耦合时包络差频导致慢拍振，而两模式仍独立守恒各自能量。
(4) 初始系统弹簧势能为 \(E(0)=\frac12(k_0+k)A^2\)，正交模式能量之和为 \(A^2m(\omega_+^2+\omega_-^2)/4\)，两式一致。若 \(k\to0\)，两频率重合，\(x_2(t)\equiv0\)；耦合增大只使反对称频率上升、对称频率不变，绝不会导致任一模式无缘无故软化。
